% TikZarium CP09 Visual Guide — XeLaTeX page template
% All interface images are included from a generated evidence manifest.
\documentclass[10pt,a4paper,landscape]{article}
\usepackage[margin=8mm,headheight=13pt,headsep=3mm,footskip=4mm,includeheadfoot]{geometry}
\usepackage{fontspec}
\usepackage{polyglossia}
\setdefaultlanguage{russian}
\setotherlanguage{english}
\setmainfont{DejaVu Sans}
\setsansfont{DejaVu Sans}
\setmonofont{DejaVu Sans Mono}[Scale=MatchLowercase]
\usepackage{xcolor}
\definecolor{Ink}{HTML}{263238}
\definecolor{Muted}{HTML}{65747A}
\definecolor{Teal}{HTML}{007C83}
\definecolor{Line}{HTML}{D5DEE0}
\definecolor{Panel}{HTML}{F3F7F7}
\usepackage{graphicx}
\usepackage{microtype}
\usepackage{booktabs}
\usepackage{enumitem}
\usepackage{fancyhdr}
\usepackage{hyperref}
\hypersetup{colorlinks=true,linkcolor=Teal,urlcolor=Teal,pdftitle={TikZarium — визуальный гид по реальному интерфейсу CP09},pdfauthor={TikZarium}}
\setlength{\parindent}{0pt}
\setlength{\parskip}{4pt}
\setlist[itemize]{leftmargin=1.2em,itemsep=1pt,topsep=2pt}
\pagestyle{fancy}
\fancyhf{}
\fancyhead[L]{\small\sffamily\color{Muted}TikZarium · CP09 · интерфейс и рабочий процесс}
\fancyhead[R]{\small\sffamily\color{Muted}Основано на наблюдаемых экранах}
\fancyfoot[L]{\small\color{Muted}CP09 · оригинальный полноэкранный PNG сохранён в assets.}
\fancyfoot[R]{\small\sffamily\color{Muted}\thepage}
\renewcommand{\headrulewidth}{0.4pt}
\renewcommand{\footrulewidth}{0pt}
\newcommand{\guideTitle}[1]{{\Large\sffamily\bfseries\color{Ink}#1}\par\vspace{3pt}\color{Teal}\rule{\linewidth}{1.1pt}\color{Ink}\vspace{5pt}}
\newcommand{\sectionTitle}[1]{{\LARGE\sffamily\bfseries\color{Ink}#1}\par\vspace{4pt}}
\newcommand{\eyebrow}[1]{{\small\sffamily\bfseries\color{Teal}\MakeUppercase{#1}}\par\vspace{2pt}}
\newcommand{\meta}[1]{{\footnotesize\sffamily\color{Muted}#1}}
\newcommand{\callout}[2]{\colorbox{Panel}{\parbox{\dimexpr\linewidth-2\fboxsep}{\textcolor{Teal}{\sffamily\bfseries #1}\hspace{0.5em}#2}}}
\newcommand{\screenImage}[2][]{\begin{center}\includegraphics[width=\linewidth,height=0.73\textheight,keepaspectratio,#1]{\detokenize{#2}}\end{center}}
\begin{document}
\thispagestyle{empty}
\vspace*{9mm}
{\sffamily\bfseries\color{Teal}\large TIKZARIUM · CP09 · VISUAL GUIDE}\par
\vspace{8mm}
{\sffamily\bfseries\color{Ink}\fontsize{30}{35}\selectfont Реальный интерфейс\par}
{\sffamily\bfseries\color{Ink}\fontsize{30}{35}\selectfont и путь от исходника до графика\par}
\vspace{6mm}
{\large\color{Muted}Практический гид по наблюдаемым экранам, действиям редактора и границам подтверждённого поведения.}\par
\vspace{10mm}
\callout{Как читать кадры}{Полноэкранный PNG включён в пакет без изменения пикселей. Подпись объясняет действие, видимый результат и его конкретную границу.}\par
\vfill
\begin{tabular}{@{}p{0.23\linewidth}p{0.69\linewidth}@{}}
\toprule
\textcolor{Teal}{\sffamily\bfseries Сборка} & \GuideBuild \\
\textcolor{Teal}{\sffamily\bfseries Исходный commit} & \GuideCommit \\
\textcolor{Teal}{\sffamily\bfseries Основное окружение} & \GuideEnvironment \\
\textcolor{Teal}{\sffamily\bfseries Текущий run} & \GuideRun \\
\bottomrule
\end{tabular}
\vspace{8mm}
\newpage
\sectionTitle{Содержание и рамка наблюдений}
\begin{minipage}[t]{0.54\linewidth}
\eyebrow{Навигация}
\GuideContents
\end{minipage}\hfill
\begin{minipage}[t]{0.40\linewidth}
\eyebrow{Что подтверждает экран}
Кадр показывает конкретный интерфейсный шаг в CP09. Полный source SHA, artifact ID/SHA и producer report связываются в \texttt{resolved\_manifest.json}.

\eyebrow{Границы операции}
Подписи относятся к видимому экрану и указанной операции. Они не доказывают поддержку другого источника, пакета или математического случая.

\eyebrow{Честность платформы}
ОС и размеры экрана берутся из producer report каждого capture. Linux/Xvfb не означает проверку Windows 10 или macOS.
\end{minipage}
\newpage
\sectionTitle{Практический маршрут и эргономика}
\begin{minipage}[t]{0.48\linewidth}
\eyebrow{Короткий сценарий}
\begin{enumerate}[leftmargin=1.3em,itemsep=4pt]
  \item Проверьте имя активного файла и видимый диапазон исходника.
  \item Выберите объект и сверьте его строку, параметр и исходный литерал.
  \item Введите значение как черновик; до Apply проверьте его ещё раз.
  \item После Apply проверьте текстовый патч и результат новой сборки.
  \item Если правка нежелательна, отмените её и сверяйте восстановление с исходным файлом.
  \item Сохраните, закройте и откройте документ заново; затем проверьте PDF и журнал.
\end{enumerate}
\end{minipage}\hfill
\begin{minipage}[t]{0.47\linewidth}
\eyebrow{Эргономика: что видно и что остаётся гипотезой}
Инспектор и редактор находятся рядом, поэтому перед изменением удобно читать источник и свойство в одном поле зрения. Видимая строка и имя активного файла снижают риск редактировать не тот документ. F1 открывает контекстную справку в зафиксированной сцене.

Клавиатурная палитра и сочетания клавиш не включены в рекомендации, если они не показаны в кадре или журнале.

\GuideReplayErgo
\end{minipage}
\newpage
% Generated screen pages follow.
\GuidePages
\end{document}
